\documentclass[10pt,a4paper]{amsart}
\usepackage{amsmath,amssymb,amsthm,graphicx}
% Retained minimal preamble for the Stacks Project chapter more-morphisms.tex (More on Morphisms).
% Verbatim excerpt of the chapter preamble: Xy-pic setup (native preamble lines 14-18) and the native
% theorem declarations with the mathematical macros (native preamble lines 154-217). Copied mechanically;
% nothing added, removed or redefined. The wrapper class and the amsmath/amssymb/amsthm/graphicx base
% packages are supplied by the assembler, so the class/externaldocument part of the native preamble is omitted.
% For commutative diagrams we use Xy-pic
\usepackage[all]{xy}

% We use 2cell for 2-commutative diagrams.
\xyoption{2cell}
% Theorem environments.
%
\theoremstyle{plain}
\newtheorem{theorem}[subsection]{Theorem}
\newtheorem{proposition}[subsection]{Proposition}
\newtheorem{lemma}[subsection]{Lemma}

\theoremstyle{definition}
\newtheorem{definition}[subsection]{Definition}
\newtheorem{example}[subsection]{Example}
\newtheorem{exercise}[subsection]{Exercise}
\newtheorem{situation}[subsection]{Situation}

\theoremstyle{remark}
\newtheorem{remark}[subsection]{Remark}
\newtheorem{remarks}[subsection]{Remarks}

\numberwithin{equation}{subsection}

% Macros
%
\def\lim{\mathop{\mathrm{lim}}\nolimits}
\def\colim{\mathop{\mathrm{colim}}\nolimits}
\def\Spec{\mathop{\mathrm{Spec}}}
\def\Hom{\mathop{\mathrm{Hom}}\nolimits}
\def\Ext{\mathop{\mathrm{Ext}}\nolimits}
\def\SheafHom{\mathop{\mathcal{H}\!\mathit{om}}\nolimits}
\def\SheafExt{\mathop{\mathcal{E}\!\mathit{xt}}\nolimits}
\def\Sch{\mathit{Sch}}
\def\Mor{\mathop{\mathrm{Mor}}\nolimits}
\def\Ob{\mathop{\mathrm{Ob}}\nolimits}
\def\Sh{\mathop{\mathit{Sh}}\nolimits}
\def\NL{\mathop{N\!L}\nolimits}
\def\CH{\mathop{\mathrm{CH}}\nolimits}
\def\proetale{{pro\text{-}\acute{e}tale}}
\def\etale{{\acute{e}tale}}
\def\QCoh{\mathit{QCoh}}
\def\Ker{\mathop{\mathrm{Ker}}}
\def\Im{\mathop{\mathrm{Im}}}
\def\Coker{\mathop{\mathrm{Coker}}}
\def\Coim{\mathop{\mathrm{Coim}}}

% Boxtimes
%
\DeclareMathSymbol{\boxtimes}{\mathbin}{AMSa}{"02}

%
% Macros for moduli stacks/spaces
%
\def\QCohstack{\mathcal{QC}\!\mathit{oh}}
\def\Cohstack{\mathcal{C}\!\mathit{oh}}
\def\Spacesstack{\mathcal{S}\!\mathit{paces}}
\def\Quotfunctor{\mathrm{Quot}}
\def\Hilbfunctor{\mathrm{Hilb}}
\def\Curvesstack{\mathcal{C}\!\mathit{urves}}
\def\Polarizedstack{\mathcal{P}\!\mathit{olarized}}
\def\Complexesstack{\mathcal{C}\!\mathit{omplexes}}
% \Pic is the operator that assigns to X its picard group, usage \Pic(X)
% \Picardstack_{X/B} denotes the Picard stack of X over B
% \Picardfunctor_{X/B} denotes the Picard functor of X over B
\def\Pic{\mathop{\mathrm{Pic}}\nolimits}
\def\Picardstack{\mathcal{P}\!\mathit{ic}}
\def\Picardfunctor{\mathrm{Pic}}
\def\Deformationcategory{\mathcal{D}\!\mathit{ef}}
\usepackage{cleveref}
\crefname{originalreference}{original reference}{original references}
\Crefname{originalreference}{Original reference}{Original references}
\makeatletter
\newlabel{algebra-lemma-compose-finite-type}{{00F4}{source}{}{sourcepointer-1}{}}
\newlabel{algebra-lemma-compose-finite-type@cref}{{[originalreference][0][]00F4}{[1][0][]source}{}{}{}}
\makeatother
\crefname{originalreference}{original reference}{original references}
\Crefname{originalreference}{Original reference}{Original references}
\makeatletter
\newlabel{algebra-lemma-finite-is-integral}{{00GK}{source}{}{sourcepointer-2}{}}
\newlabel{algebra-lemma-finite-is-integral@cref}{{[originalreference][0][]00GK}{[1][0][]source}{}{}{}}
\makeatother
\title{A universal ring for idempotents of a finite finitely presented algebra}
\author{The Stacks Project Authors}
\date{stacks-project:more-morphisms.tex:local-snapshot-20261011; GNU FDL 1.2 or later}
\begin{document}
\maketitle
\noindent\textit{Editorial note.} Counted span is the author's own proof environment, source lines 22799-22832, immediately following the statement at 22790-22797.\par\medskip
\noindent\textit{Editorial note.} Two cross-chapter pointers (Algebra) are recorded in assembly.external\_references.\par\medskip
\noindent\textit{Editorial note.} The universal ring construction is carried out explicitly with generators and equations inside the span.\par\medskip
\noindent\textit{Allowed standard prerequisites (editorial).}\begin{itemize}
\item Finiteness of an integral element. Algebra, lemma algebra-lemma-finite-is-integral, cited at proof line 22802. Produces a monic polynomial relation for each chosen generator of B over A.
\item Composition of finite type ring maps. Algebra, lemma algebra-lemma-compose-finite-type, cited in the proof. Keeps the constructions of finite presentation while adjoining the universal idempotent data.
\item Finite presentation of a quotient of a finite free algebra. used at proof lines 22803 onward. Presents B as a quotient of a finite free A-algebra, so that idempotents can be handled by polynomial equations with universal coefficients.
\end{itemize}\medskip
\ifdefined\hypertarget\hypertarget{sourcepointer-1}{}\fi
\noindent\textit{Source reference.} Original reference 00F4: focused-seeds/geometry/S9-morph-open-closed-in-finite/ref-tags.txt, source lines 4--4. Navigation pointer only: the counted proof cites this statement by \textbackslash{}ref as a standing input; the deduction the counted proof has to make is carried out inside the counted span, and no part of that deduction is delegated to this pointer.\par\medskip
\ifdefined\hypertarget\hypertarget{sourcepointer-2}{}\fi
\noindent\textit{Source reference.} Original reference 00GK: focused-seeds/geometry/S9-morph-open-closed-in-finite/ref-tags.txt, source lines 5--5. Navigation pointer only: the counted proof cites this statement by \textbackslash{}ref as a standing input; the deduction the counted proof has to make is carried out inside the counted span, and no part of that deduction is delegated to this pointer.\par\medskip
\begin{lemma}
\label{lemma-open-and-closed-in-finite}
Let $A \to B$ be a ring map which is finite and of finite presentation.
There exists a finitely presented ring map $A \to A_{univ}$
and an idempotent $e_{univ} \in B \otimes_A A_{univ}$
such that for any ring map $A \to A'$ and idempotent $e \in B \otimes_A A'$
there is a ring map $A_{univ} \to A'$ mapping $e_{univ}$ to $e$.
\end{lemma}

\begin{proof}
Choose $b_1, \ldots, b_n \in B$ generating $B$ as an $A$-module.
For each $i$ choose a monic $P_i \in A[x]$ such that $P_i(b_i) = 0$
in $B$, see Algebra, Lemma \ref{algebra-lemma-finite-is-integral}.
Thus $B$ is a quotient of the finite free $A$-algebra
$B' = A[x_1, \ldots, x_n]/(P_1(x_1), \ldots, P_n(x_n))$.
Let $J \subset B'$ be the kernel of the surjection $B' \to B$.
Then $J =(f_1, \ldots, f_m)$ is finitely generated as $B$
is a finitely generated $A$-algebra, see
Algebra, Lemma \ref{algebra-lemma-compose-finite-type}.
Choose an $A$-basis $b'_1, \ldots, b'_N$ of $B'$.
Consider the algebra
$$
A_{univ} = A[z_1, \ldots, z_N, y_1, \ldots, y_m]/I
$$
where $I$ is the ideal generated by the coefficients in
$A[z_1, \ldots, z_n, y_1, \ldots, y_m]$
of the basis elements $b'_1, \ldots, b'_N$ of the expression
$$
(\sum z_j b'_j)^2 - \sum z_j b'_j + \sum y_k f_k
$$
in $B'[z_1, \ldots, z_N, y_1, \ldots, y_m]$. By construction
the element $\sum z_j b'_j$ maps to an idempotent $e_{univ}$ in the
algebra $B \otimes_A A_{univ}$. Moreover, if $e \in B \otimes_A A'$
is an idempotent, then we can lift $e$ to an element of the form
$\sum b'_j \otimes a'_j$ in $B' \otimes_A A'$ and we can find
$a''_k \in A'$ such that
$$
(\sum b'_j \otimes a'_j)^2 - \sum b'_j \otimes a'_j + \sum f_k \otimes a''_k
$$
is zero in $B' \otimes_A A'$. Hence we get an $A$-algebra map
$A_{univ} \to A$ sending $z_j$ to $a'_j$ and $y_k$ to $a''_k$
mapping $e_{univ}$ to $e$. This finishes the proof.
\end{proof}

\end{document}
